\documentclass{article}

\usepackage[nonatbib]{neurips_2026}

\usepackage[utf8]{inputenc}
\usepackage[T1]{fontenc}
\usepackage{hyperref}
\usepackage{url}
\usepackage{booktabs}
\usepackage{amsfonts}
\usepackage{nicefrac}
\usepackage{microtype}
\usepackage{xcolor}
\usepackage{graphicx}
\usepackage{amsmath}

\title{SMC-QA: Selective Memory Consolidation with Query-Aware Retrieval for Long-Context Question Answering\\
\large Extended Improvement Study}

\author{
  Zheng Yifan \\
  SID: 12532245 \\
  \texttt{yifan.zheng@example.edu} \\
  \And
  Chen Zeming \\
  SID: 12531410 \\
  \texttt{zeming.chen@example.edu} \\
}

\begin{document}

\maketitle

\begin{abstract}
Long-context question answering requires systems to manage and retrieve information from extensive interaction histories. A flat memory approach treats all past information equally, which may introduce noise and miss the distinction between recent and consolidated knowledge. In this project, we build on SMC-QA, a lightweight hierarchical memory framework that organizes information into short-term memory (STM) and long-term memory (LTM). The base system uses a selective promotion strategy based on relevance, reuse frequency, and information diversity, together with rule-based query-aware retrieval routing. We evaluate SMC-QA and five baselines on the LongMemEval benchmark using both oracle and full-haystack settings, then add improvement experiments without modifying the shared project source code. The improved retrieval-depth variant increases oracle Hit@6 from 0.778 to 0.793 and Recall@6 from 0.612 to 0.631. On the harder S-Cleaned setting, a relevance-heavy tuned configuration improves SMC-QA from 0.420 to 0.510 Hit@6, outperforming Flat Memory on the 100-instance evaluation subset. We also prepare an optional LLM answer generation pipeline, although no API-based generation result is reported because the environment did not contain an OpenAI API key.
\end{abstract}

%======================================================================
\section{Introduction}
%======================================================================

Long-context question answering (QA) is a challenging task that requires systems to locate relevant evidence within large volumes of historical interactions. Unlike standard QA, where the context is typically short and self-contained, long-context QA involves multi-session dialogues, knowledge updates, and temporal reasoning across extended interaction histories.

A fundamental challenge in this setting is \textbf{memory management}: how should a system decide what information to retain, what to discard, and where to search when a new question arrives? A simple approach is to store all past information in a single pool and retrieve from it directly, but this leads to increased noise as irrelevant content accumulates.

Inspired by cognitive science concepts of short-term and long-term memory, as well as recent work on hierarchical memory systems for LLMs~\cite{himem, rmm, lightmem}, we propose \textbf{SMC-QA} (Selective Memory Consolidation for Question Answering). Our system:

\begin{itemize}
    \item Organizes information into a \textbf{short-term memory} (STM) for recent content and a \textbf{long-term memory} (LTM) for consolidated knowledge.
    \item Uses a lightweight \textbf{promotion score} combining relevance, reuse frequency, and diversity to decide which STM items should be promoted to LTM.
    \item Applies \textbf{rule-based query-aware routing} to decide whether to search STM, LTM, or both based on query characteristics.
\end{itemize}

We compare SMC-QA against five baselines on the LongMemEval benchmark~\cite{longmemeval}, conduct ablation studies to understand the contribution of each component, and further evaluate whether small retrieval and promotion changes can address the observed weaknesses of the base configuration.

%======================================================================
\section{Related Work}
%======================================================================

\textbf{LongMemEval}~\cite{longmemeval} provides a benchmark for evaluating long-term interactive memory across five capabilities: information extraction, multi-session reasoning, knowledge updates, temporal reasoning, and abstention. We use this benchmark as our primary evaluation dataset.

\textbf{HiMem}~\cite{himem} introduces hierarchical long-term memory for long-horizon dialogue agents, supporting both hierarchical memory construction and hybrid retrieval. Our STM/LTM structure draws inspiration from this work but uses a simpler, non-parametric design.

\textbf{Reflective Memory Management (RMM)}~\cite{rmm} proposes multi-granularity memory summarization across utterances, turns, and sessions. While we do not implement summary-based promotion in our current system, this work motivates our consideration of memory consolidation.

\textbf{LightMem}~\cite{lightmem} modularizes memory into retrieval, writing, and long-term consolidation components. Our system adopts a similar modular philosophy but focuses specifically on the promotion decision and retrieval routing.

%======================================================================
\section{Method}
%======================================================================

\subsection{System Overview}

SMC-QA processes a stream of text chunks from historical interactions and organizes them into a two-level memory hierarchy. When a query arrives, the system routes it to the appropriate memory level(s) and retrieves the most relevant evidence.

\subsection{Memory Structure}

\textbf{Short-Term Memory (STM)} operates as a fixed-capacity FIFO buffer (capacity $C_s = 40$ chunks). New chunks enter STM, and when capacity is exceeded, the oldest chunks are evicted. STM preserves recent, detailed information.

\textbf{Long-Term Memory (LTM)} stores promoted items with a maximum capacity of $C_l = 120$ chunks. Items enter LTM only through the promotion process. When LTM exceeds capacity, the least-accessed items are evicted. A duplicate detection mechanism (cosine similarity $\geq 0.85$) prevents redundant entries.

\subsection{Selective Memory Consolidation}

Every $f = 5$ chunks, a consolidation step evaluates STM items for promotion to LTM. Each candidate receives a composite score:

\begin{equation}
    \text{score} = w_r \cdot \text{relevance} + w_u \cdot \text{reuse} + w_d \cdot \text{diversity}
\end{equation}

where $w_r = 0.4$, $w_u = 0.4$, $w_d = 0.2$.

\textbf{Relevance} measures the average cosine similarity between the chunk embedding and the embeddings of recent queries (window of 5).

\textbf{Reuse} captures how frequently the chunk has been retrieved, normalized by the maximum hit count in STM.

\textbf{Diversity} measures novelty relative to existing LTM content: $\text{diversity} = 1 - \max_{l \in \text{LTM}} \text{sim}(c, l)$.

Only the top-$k$ candidates ($k=3$) with scores above a threshold ($\tau = 0.35$) are promoted.

\subsection{Query-Aware Retrieval Routing}

When a query arrives, the system applies three rules in sequence:

\begin{enumerate}
    \item \textbf{Rule 1 (LTM-first)}: If the query contains temporal keywords (e.g., ``before'', ``earlier'', ``previously'', ``first''), route to LTM.
    \item \textbf{Rule 2 (STM-first)}: If the query's average similarity to recent STM chunks exceeds its similarity to LTM by a margin $\delta = 0.05$, route to STM.
    \item \textbf{Rule 3 (Hybrid)}: Otherwise, retrieve from both STM and LTM, merge results, and re-rank.
\end{enumerate}

For single-path retrieval, we return top-4 results. For hybrid retrieval, we take top-4 from each level, merge, deduplicate, and re-rank to return top-6.

\subsection{Improved Retrieval-Depth Variant}

During the improvement study, we found a mismatch between the retrieval setting and the evaluation metric: single-path STM or LTM routing returns only four chunks, while the main metric is Hit@6 / Recall@6. We therefore add a minimal improved variant that keeps the base memory construction and routing logic but widens the per-memory retrieval depth.

For the oracle setting, the improved variant uses the base promotion weights $(w_r, w_u, w_d)=(0.4,0.4,0.2)$, a slightly stricter promotion threshold $\tau=0.40$, top-3 promotion, and top-8 retrieval from STM and LTM before final top-6 re-ranking. For S-Cleaned, we additionally test a relevance-heavy promotion score $(0.5,0.3,0.2)$ with $\tau=0.20$, because the full-haystack setting contains substantially more distractor content and benefits from promoting semantically relevant chunks earlier. These variants are implemented as separate scripts under \texttt{Improvement/}; the shared project source files are left unchanged.

%======================================================================
\section{Experimental Setup}
%======================================================================

\subsection{Dataset}

We use the \textbf{LongMemEval} benchmark~\cite{longmemeval}, which contains 500 evaluation instances across six question types: single-session-user, single-session-assistant, single-session-preference, multi-session, temporal-reasoning, and knowledge-update.

We evaluate in two settings:
\begin{itemize}
    \item \textbf{Oracle}: Each instance includes only answer-relevant sessions ($\sim$28 turns/instance).
    \item \textbf{S-Cleaned (Full Haystack)}: Each instance includes $\sim$50 sessions with $\sim$500 turns, where answer-relevant content is mixed with distractors. We truncate to 200 turns per instance for computational tractability.
\end{itemize}

For oracle experiments, we sample 200 instances (50 dev / 150 test) with seed 42 and run formal experiments with three seeds (42, 43, 44).

\subsection{Baselines}

\begin{enumerate}
    \item \textbf{Flat Memory}: All chunks in one pool; retrieve top-6 by cosine similarity.
    \item \textbf{STM-only}: Retrieve only from the most recent 40 chunks.
    \item \textbf{LTM-only}: Build LTM via our promotion method; retrieve only from LTM.
    \item \textbf{Naive STM+LTM}: Build STM/LTM; always retrieve from both (no routing).
    \item \textbf{Frequency-based Promotion}: Promote to LTM based on retrieval frequency only.
\end{enumerate}

\subsection{Metrics}

\begin{itemize}
    \item \textbf{Hit@6}: Whether at least one gold evidence turn is found in the top-6 retrieved chunks.
    \item \textbf{Recall@6}: Fraction of gold evidence turns covered by the top-6 retrieved chunks.
    \item \textbf{Contains}: Whether the concatenated top-3 retrieved chunks contain the gold answer string.
\end{itemize}

We match retrieved chunks to gold evidence using word-level overlap with a threshold of 50\%.

\subsection{Implementation}

We use \texttt{all-MiniLM-L6-v2} from sentence-transformers for embedding (384-dimensional). All retrieval uses brute-force cosine similarity. Text is chunked at 120 words with 30-word overlap, preserving natural turn boundaries when turns are shorter than 160 words.

\subsection{Improvement Protocol}

Baseline run artifacts are stored under \texttt{Note/}, while improvement artifacts are stored under \texttt{Improvement/}. This separation keeps the shared project implementation intact while making the extended experiments easy to audit. The main improvement scripts are:

\begin{itemize}
    \item \texttt{tune\_smc.py}: promotion-only tuning on the oracle split.
    \item \texttt{tune\_smc\_retrieval.py}: joint promotion and retrieval-depth tuning.
    \item \texttt{supplementary\_experiments.py}: question-type breakdown, LTM quality analysis, and case studies.
    \item \texttt{tune\_s\_cleaned.py}: S-Cleaned-specific promotion and retrieval tuning.
    \item \texttt{answer\_generation.py}: optional answer generation using the OpenAI Responses API.
\end{itemize}

The S-Cleaned tuning experiment uses the same 100-instance subset as the baseline run. We tune on the first 30 instances and report the final comparison on all 100 instances, each truncated to 200 turns. The answer generation script was executed in dry-run mode only, producing prompts and metadata for 20 seed-42 test instances, because \texttt{OPENAI\_API\_KEY} was not set in the execution environment.

The experiments reuse the conda environment \texttt{ml\_project\_repro}. The core dependencies are installed with \texttt{python -m pip install sentence-transformers tqdm numpy}, and the exact package snapshot is saved to \texttt{Improvement/results/pip\_freeze.txt}. All improvement runs are executed through wrapper scripts or in-process parameter overrides; no changes are written back to \texttt{src/}, \texttt{run\_experiments.py}, \texttt{run\_s\_cleaned.py}, or \texttt{run\_case\_analysis.py}.

%======================================================================
\section{Results}
%======================================================================

\subsection{Oracle Setting}

Table~\ref{tab:main} shows results on the oracle dataset (150 test instances, mean $\pm$ std across 3 seeds).

\begin{table}[h]
\centering
\caption{Main results on LongMemEval Oracle (150 test instances, 3 seeds).}
\label{tab:main}
\begin{tabular}{lccc}
\toprule
\textbf{Method} & \textbf{Hit@6} & \textbf{Recall@6} & \textbf{Contains} \\
\midrule
STM-only & 0.642 $\pm$ 0.003 & 0.478 $\pm$ 0.002 & 0.342 $\pm$ 0.042 \\
LTM-only & 0.713 $\pm$ 0.019 & 0.549 $\pm$ 0.014 & 0.309 $\pm$ 0.044 \\
Flat Memory & 0.782 $\pm$ 0.019 & 0.623 $\pm$ 0.021 & 0.298 $\pm$ 0.049 \\
Naive STM+LTM & 0.776 $\pm$ 0.017 & 0.614 $\pm$ 0.018 & 0.329 $\pm$ 0.039 \\
Freq Promotion & 0.804 $\pm$ 0.013 & 0.641 $\pm$ 0.021 & 0.327 $\pm$ 0.045 \\
SMC-QA (Ours) & 0.778 $\pm$ 0.022 & 0.612 $\pm$ 0.018 & \textbf{0.333 $\pm$ 0.030} \\
\bottomrule
\end{tabular}
\end{table}

\textbf{Key observations:}
\begin{itemize}
    \item \textbf{Hierarchical memory helps}: STM-only (0.642 Hit@6) significantly underperforms all methods that include LTM, confirming that relying solely on recent memory is insufficient.
    \item \textbf{LTM adds value}: LTM-only (0.713) outperforms STM-only by 7.1\%, showing that promoted content provides useful signal even without short-term context.
    \item \textbf{Promotion strategies are competitive}: Both Freq Promotion (0.804) and SMC-QA (0.778) perform comparably to or better than Flat Memory (0.782) on Hit@6, despite having access to only a subset of chunks through their memory structures.
    \item \textbf{Contains metric}: SMC-QA achieves the highest Contains score (0.333) with the lowest variance (0.030), suggesting it retrieves more answer-relevant content.
\end{itemize}

\subsection{Promotion-Only Tuning}

Before changing retrieval depth, we first tune only the promotion weights, promotion threshold, promotion top-$k$, and routing mode. The search covers baseline, relevance-heavy, relevance-strong, and diversity-heavy weighting schemes; thresholds $\{0.25,0.30,0.35,0.40\}$; promotion top-$k \in \{3,5\}$; and auto versus hybrid routing. Table~\ref{tab:promotion-tuning} lists the best test configurations from this first stage.

\begin{table}[h]
\centering
\caption{Best promotion-only tuning results on the oracle test split.}
\label{tab:promotion-tuning}
\resizebox{\textwidth}{!}{
\begin{tabular}{lccc}
\toprule
\textbf{Config} & \textbf{Hit@6} & \textbf{Recall@6} & \textbf{Contains} \\
\midrule
\texttt{baseline\_thr0.40\_top3\_auto} & 0.776 $\pm$ 0.023 & 0.614 $\pm$ 0.016 & \textbf{0.336 $\pm$ 0.035} \\
\texttt{baseline\_thr0.25\_top3\_auto} & 0.776 $\pm$ 0.017 & 0.609 $\pm$ 0.013 & 0.333 $\pm$ 0.030 \\
\texttt{div\_heavy\_thr0.40\_top5\_auto} & 0.760 $\pm$ 0.016 & 0.593 $\pm$ 0.016 & 0.333 $\pm$ 0.036 \\
\bottomrule
\end{tabular}
}
\end{table}

The first-stage conclusion is negative but useful: promotion-only tuning slightly improves Contains, but it does not reliably improve Hit@6 or Recall@6 and still fails to outperform the frequency-promotion baseline. This motivates the second-stage focus on retrieval top-$k$.

\subsection{Improved Oracle Retrieval Depth}

Table~\ref{tab:retrieval-tuned} compares the base SMC-QA against the best retrieval-depth variant. Increasing the STM/LTM candidate depth from four to eight improves both Hit@6 and Recall@6 while leaving Contains nearly unchanged. The tuned model still does not surpass frequency-based promotion on Hit@6, but it closes part of the gap while preserving the query-aware SMC-QA structure.

The retrieval top-$k$ search first selects candidate configurations on the 50-instance dev set, then evaluates the best configurations across three test seeds. The strongest two dev configurations both reach 0.800 Hit@6 and 0.643 Recall@6, showing that the main gain comes from retrieving at least six to eight candidates per memory level.

\begin{table}[h]
\centering
\caption{Retrieval top-$k$ tuning details from the improvement run log.}
\label{tab:retrieval-search}
\resizebox{\textwidth}{!}{
\begin{tabular}{lccc}
\toprule
\textbf{Config} & \textbf{Hit@6} & \textbf{Recall@6} & \textbf{Contains} \\
\midrule
\texttt{baseline\_thr0.40\_top3\_auto\_stm6\_ltm6} & 0.789 $\pm$ 0.033 & 0.628 $\pm$ 0.024 & 0.329 $\pm$ 0.033 \\
\texttt{baseline\_thr0.40\_top3\_auto\_stm8\_ltm8} & \textbf{0.793 $\pm$ 0.030} & \textbf{0.631 $\pm$ 0.024} & 0.329 $\pm$ 0.033 \\
\texttt{div\_heavy\_thr0.40\_top5\_auto\_stm6\_ltm4} & 0.769 $\pm$ 0.019 & 0.606 $\pm$ 0.017 & 0.327 $\pm$ 0.033 \\
\bottomrule
\end{tabular}
}
\end{table}

\begin{table}[h]
\centering
\caption{Oracle improvement from retrieval-depth tuning (150 test instances, 3 seeds).}
\label{tab:retrieval-tuned}
\begin{tabular}{lccc}
\toprule
\textbf{Method} & \textbf{Hit@6} & \textbf{Recall@6} & \textbf{Contains} \\
\midrule
Base SMC-QA & 0.778 $\pm$ 0.022 & 0.612 $\pm$ 0.018 & \textbf{0.333 $\pm$ 0.030} \\
Retrieval-tuned SMC-QA & 0.793 $\pm$ 0.030 & 0.631 $\pm$ 0.024 & 0.329 $\pm$ 0.033 \\
Freq Promotion & \textbf{0.804 $\pm$ 0.013} & \textbf{0.641 $\pm$ 0.021} & 0.327 $\pm$ 0.045 \\
\bottomrule
\end{tabular}
\end{table}

The question-type breakdown in Table~\ref{tab:question-type} shows that the tuned variant mainly helps \textit{temporal-reasoning} and \textit{single-session-assistant} questions. However, it regresses on \textit{knowledge-update}, suggesting that wider retrieval can introduce older or conflicting evidence when the question asks for the most recent state.

\begin{table}[h]
\centering
\caption{Question-type comparison between base and retrieval-tuned SMC-QA on seed 42.}
\label{tab:question-type}
\begin{tabular}{lcccc}
\toprule
\textbf{Question Type} & \textbf{Base Hit} & \textbf{Tuned Hit} & \textbf{Base Recall} & \textbf{Tuned Recall} \\
\midrule
single-session-assistant & 0.588 & \textbf{0.647} & 0.588 & \textbf{0.647} \\
temporal-reasoning & 0.769 & \textbf{0.821} & 0.544 & \textbf{0.569} \\
multi-session & \textbf{0.750} & 0.722 & 0.485 & \textbf{0.501} \\
knowledge-update & \textbf{0.778} & 0.741 & \textbf{0.580} & 0.543 \\
\bottomrule
\end{tabular}
\end{table}

\subsection{Full Haystack Setting (S-Cleaned)}

Table~\ref{tab:scleaned} shows results when answer-relevant content is mixed with $\sim$200 distractor turns. The baseline run showed that Flat Memory outperformed SMC-QA in this setting. After tuning the S-Cleaned-specific promotion threshold, promotion weights, and retrieval depth, the best SMC-QA variant improves substantially and becomes the strongest method on this 100-instance subset.

\begin{table}[h]
\centering
\caption{Results on LongMemEval S-Cleaned before and after tuning (100 instances, 200 turns each).}
\label{tab:scleaned}
\begin{tabular}{lccc}
\toprule
\textbf{Method} & \textbf{Hit@6} & \textbf{Recall@6} & \textbf{Contains} \\
\midrule
Flat Memory & 0.460 & 0.315 & 0.260 \\
Freq Promotion & 0.430 & 0.292 & 0.250 \\
Base SMC-QA & 0.420 & 0.290 & 0.250 \\
Tuned SMC-QA & \textbf{0.510} & \textbf{0.360} & \textbf{0.270} \\
\bottomrule
\end{tabular}
\end{table}

The best S-Cleaned configuration is \texttt{rel\_heavy\_thr0.20\_top3\_stm8\_ltm8}. Compared with base SMC-QA, it improves Hit@6 by 9.0 percentage points and Recall@6 by 7.0 percentage points. This result indicates that the base hyperparameters were too conservative for full-haystack retrieval: lowering the promotion threshold and increasing retrieval depth helps the model preserve and expose relevant evidence among distractors.

\subsection{Ablation Study}

Table~\ref{tab:ablation} shows the effect of removing each component from SMC-QA.

\begin{table}[h]
\centering
\caption{Ablation results on LongMemEval Oracle (150 test instances).}
\label{tab:ablation}
\begin{tabular}{lccc}
\toprule
\textbf{Configuration} & \textbf{Hit@6} & \textbf{Recall@6} & \textbf{Contains} \\
\midrule
SMC-QA (full) & 0.753 & 0.588 & 0.327 \\
$-$ relevance & 0.760 & 0.589 & 0.333 \\
$-$ diversity & 0.740 & 0.570 & 0.313 \\
$-$ routing & 0.753 & 0.589 & 0.320 \\
\bottomrule
\end{tabular}
\end{table}

\textbf{Diversity matters most}: Removing the diversity component causes the largest drop ($-$1.3\% Hit@6, $-$1.8\% Recall@6), indicating that preventing duplicate information from entering LTM is the most valuable aspect of our promotion strategy.

\textbf{Relevance has minimal impact}: Removing relevance slightly improves performance, suggesting that in the oracle setting where most content is already relevant, this signal provides less discriminative power.

\textbf{Routing has limited effect}: Removing query-aware routing shows no change in Hit@6 and slight decrease in Contains, suggesting that the rule-based routing provides modest benefits in this dataset.

\subsection{LTM Quality Analysis}

To understand why retrieval tuning helps, we also evaluate the quality of the constructed LTM itself. Table~\ref{tab:ltm-quality} reports the average LTM size, the fraction of LTM items that contain the answer string, and whether gold evidence is retrievable directly from LTM.

\begin{table}[h]
\centering
\caption{LTM quality analysis on seed 42.}
\label{tab:ltm-quality}
\begin{tabular}{lcccc}
\toprule
\textbf{Method} & \textbf{LTM Size} & \textbf{Answer Item Rate} & \textbf{Gold Hit} & \textbf{Gold Recall} \\
\midrule
Freq Promotion & 8.7 & \textbf{0.144} & 0.787 & 0.607 \\
Base SMC-QA & \textbf{13.6} & 0.100 & \textbf{0.800} & \textbf{0.650} \\
Retrieval-tuned SMC-QA & 10.1 & 0.120 & 0.780 & 0.603 \\
\bottomrule
\end{tabular}
\end{table}

The tuned variant does not improve intrinsic LTM quality over base SMC-QA. Its final retrieval gain therefore comes mainly from exposing more STM/LTM candidates before the final top-6 ranking, rather than from building a better LTM. This also explains why frequency-based promotion remains a strong baseline: although its LTM is smaller, a larger fraction of its LTM items directly contain answer strings.

\subsection{Optional Answer Generation Preparation}

We also prepared a downstream answer generation experiment to test whether retrieval improvements translate into final natural-language answers. The script \texttt{Improvement/answer\_generation.py} builds prompts for Flat/Frequency/SMC-style retrieved evidence and can call the OpenAI Responses API when \texttt{OPENAI\_API\_KEY} is set. In the current execution environment, no API key was available, so the experiment was run only in dry-run mode on 20 seed-42 test instances. Therefore, this report does not claim answer-generation accuracy, but the prompt files and execution metadata are available under \texttt{Improvement/results/answer\_generation\_dry\_run.json}.

\subsection{Case Analysis}

We analyze specific instances where the retrieval-tuned SMC-QA improves over or regresses from the base SMC-QA. Table~\ref{tab:case-study} summarizes representative cases from the seed-42 supplementary experiment.

\begin{table}[h]
\centering
\caption{Representative case studies from the retrieval-tuned SMC-QA analysis.}
\label{tab:case-study}
\resizebox{\textwidth}{!}{
\begin{tabular}{llcc}
\toprule
\textbf{Type} & \textbf{Case Description} & \textbf{Original Recall} & \textbf{Tuned Recall} \\
\midrule
single-session-assistant & Catalonia singer-songwriter question & 0.00 & \textbf{1.00} \\
multi-session & Initial tomatoes and cucumbers plants & 0.50 & \textbf{1.00} \\
temporal-reasoning & Charity gala versus charity bake sale & 0.00 & \textbf{0.50} \\
knowledge-update & Old sneakers location & \textbf{1.00} & 0.50 \\
knowledge-update & Alex from Germany meetups & \textbf{0.50} & 0.00 \\
multi-session & Car wash and parking ticket spending & \textbf{1.00} & 0.50 \\
\bottomrule
\end{tabular}
}
\end{table}

\textbf{Success cases} often require exposing an evidence chunk that was just outside the base top-4 candidate set. Wider candidate retrieval gives the final top-6 ranker a chance to recover these items. \textbf{Regression cases} often involve updates, locations, or spending details where older semantically similar evidence competes with the current answer. This supports the earlier conclusion that retrieval depth should be paired with recency-aware conflict handling.

%======================================================================
\section{Discussion}
%======================================================================

\textbf{When does hierarchical memory help?} Our results show that the benefit of hierarchical memory is most clear when (1) the conversation history is long enough that STM alone cannot cover relevant information, and (2) the question specifically requires long-term knowledge. In the oracle setting with limited context, the advantage is modest, but the S-Cleaned experiment shows that tuned hierarchical memory can outperform Flat Memory even when many distractors are present.

\textbf{Retrieval depth matters}: The strongest oracle improvement comes from a small mismatch in the base retrieval setup: single-path routing returns top-4 results, while evaluation is based on top-6. Expanding STM/LTM candidate retrieval to top-8 improves Hit@6 and Recall@6 without changing the base memory architecture. This suggests that SMC-QA's memory construction is useful, but the retriever needs enough candidate bandwidth to expose relevant evidence.

\textbf{Promotion strategy trade-offs}: The frequency-based baseline performs surprisingly well, suggesting that simple heuristics can be effective. Our multi-factor promotion score (relevance + reuse + diversity) shows its strength primarily through the diversity component in oracle ablations, while the S-Cleaned tuning experiment shows that relevance should receive more weight when distractors dominate. The best S-Cleaned configuration uses lower promotion threshold and wider retrieval, indicating that conservative consolidation can discard useful evidence too early.

\textbf{Remaining failure modes}: The question-type breakdown shows a clear regression on knowledge-update questions after retrieval-depth tuning. This likely happens because wider retrieval admits stale but semantically similar evidence. A better variant should add recency-aware conflict resolution, especially for questions containing cues such as ``most recent'', ``latest'', or ``updated''.

\textbf{Limitations}: (1) Our routing rules are hand-crafted and may not generalize to all query types. A learned router could be more adaptive. (2) We do not perform summary-based promotion, which could compress information more effectively. (3) The embedding model (MiniLM-L6-v2) is relatively small; stronger encoders may improve retrieval quality. (4) The improvement experiments tune on relatively small subsets, especially S-Cleaned. (5) We evaluate primarily on retrieval metrics; the answer generation experiment was prepared but not executed with an LLM API because no API key was available.

\textbf{Future work}: Extending this framework with a learned routing module, recency-aware update handling, summary-based consolidation, stronger embedding models, and answer generation evaluation would provide a more complete picture of hierarchical memory in long-context QA.

%======================================================================
\section{Conclusion}
%======================================================================

We presented an extended improvement study of SMC-QA, a lightweight hierarchical memory framework for long-context question answering. The baseline experiments show that STM+LTM memory is competitive with flat retrieval and that diversity is the most important component of the base promotion score. Our improvement experiments further show that retrieval depth is a practical bottleneck: expanding STM/LTM candidate retrieval improves oracle Hit@6 from 0.778 to 0.793 and Recall@6 from 0.612 to 0.631. On the harder S-Cleaned setting, relevance-heavy promotion with lower threshold and wider retrieval improves SMC-QA from 0.420 to 0.510 Hit@6, outperforming Flat Memory on the evaluated subset. These results suggest that hierarchical memory is promising, but its effectiveness depends strongly on promotion aggressiveness, retrieval bandwidth, and special handling of knowledge updates.

%======================================================================
\section{Contributions}
%======================================================================

\textbf{Zheng Yifan} (50\%): System design, implementation of memory structures, promotion strategies, routing module, experiment pipeline, and report writing.

\textbf{Chen Zeming} (50\%): Data processing, baseline implementation, evaluation metrics, ablation experiments, and case analysis.

%======================================================================
% References
%======================================================================

\begin{thebibliography}{4}

\bibitem{longmemeval}
Wu, D., Wang, H., Yu, W., Zhang, Y., Chang, K.-W., and Yu, D.
\newblock LongMemEval: Benchmarking Chat Assistants on Long-Term Interactive Memory.
\newblock \emph{arXiv preprint arXiv:2410.10813}, 2024.

\bibitem{himem}
Zhang, N., Yang, X., Tan, Z., Deng, W., and Wang, W.
\newblock HiMem: Hierarchical Long-Term Memory for LLM Long-Horizon Agents.
\newblock \emph{arXiv preprint arXiv:2601.06377}, 2026.

\bibitem{rmm}
Tan, Z., Yan, J., Hsu, I.-H., Han, R., Wang, Z., Le, L., Song, Y., Chen, Y., Palangi, H., Lee, G., Iyer, A.~R., Chen, T., Liu, H., Lee, C.-Y., and Pfister, T.
\newblock In Prospect and Retrospect: Reflective Memory Management for Long-term Personalized Dialogue Agents.
\newblock In \emph{Proceedings of ACL}, pages 8416--8439, 2025.

\bibitem{lightmem}
Fang, J., Deng, X., Xu, H., Jiang, Z., Tang, Y., Xu, Z., Deng, S., Yao, Y., Wang, M., Qiao, S., Chen, H., and Zhang, N.
\newblock LightMem: Lightweight and Efficient Memory-Augmented Generation.
\newblock \emph{arXiv preprint arXiv:2510.18866}, 2025.

\end{thebibliography}

\end{document}
